\begin{frame}{Définition}

    \begin{center}
        {\small  \textit{L’œuvre « profilaire » (Monjour et Vitali-Rosati, 2017) est constituée d’une page-profil, de \textit{posts}, commentaires, \textit{loves} et \textit{likes}, de demandes en amitié envoyées et acceptées, mais aussi de mouvements de \textit{scroll}, arrêts sur image, clics et autres traces que l’auteur laisse sur la plateforme sans les conscientiser comme des activités d’écriture. Chaque profil se caractérise ainsi par une face déclarée et une face automatiquement calculée qui, comme l’a formulé Louise Merzeau (2016), le suit comme une ombre.}}\\
        \textbf{Appiotti, S. et Saemmer, A. (2022), \textit{Le profil de fiction sur les réseaux sociaux Au prisme de nos lectures}}
    \end{center}
\end{frame}

\begin{frame}{Le "pacte autobiographique"}
    \begin{center}
        {\small \textit{C'est l'engagement que prend un auteur de raconter directement sa vie (ou une partie, ou un aspect de sa vie) dans un esprit de vérité. Le pacte autobiographique s'oppose au pacte de fiction. Quelqu'un qui vous propose un roman (même s'il est inspiré de sa vie) ne vous demande pas de croire pour de bon à ce qu'il raconte : mais simplement de jouer à y croire. L'autobiographe, lui, vous promet que ce que qu'il va vous dire est vrai, ou, du moins, est ce qu'il croit vrai.}}\\
        \textbf{Philippe Lejeune, \textit{Le Pacte autobiographique}, 1975.}
    \end{center}
    \textrightarrow{} Que devient le pacte autobiographique à l'heure des écritures numériques?
\end{frame}
\begin{frame}{Problématiques}
    \begin{itemize}
        \item Enjeux de la construction identitaire en ligne, et sa relation avec notre identité "réelle"
        \item Comment l'objet profil redessine-t-il la notion d'auteur et de personnage?
        \item Quelle participation pour le lecteur?
        \item Les plateformes sociales présentent-elles des contraintes éditoriales?
        \item Quelle légitimation?
    \end{itemize}
\end{frame}
